% ============================================================
%   classicthesis-config.tex
%   Configuración del paquete classicthesis
% ============================================================

\PassOptionsToPackage{eulerchapternumbers,listings,pdfspacing,subfig,beramono,eulermath,dottedtoc}{classicthesis}

% ------------------------------------------------------------
%   INFORMACIÓN DEL DOCUMENTO
% ------------------------------------------------------------

\newcommand{\myTitle}{VICON: Sistema de Visión configurable\xspace}
\newcommand{\mySubtitle}{\xspace}
\newcommand{\myName}{Carlos Manuel Gómez Jiménez\xspace}
\newcommand{\mySupervisor}{Martín González García\xspace}
\newcommand{\myDepartment}{Departamento de Ingeniería Informática y de Sistemas\xspace}
\newcommand{\myUni}{Universidad de Málaga\xspace}
\newcommand{\myLocation}{Málaga\xspace}
\newcommand{\myTime}{Septiembre 2026\xspace}
\newcommand{\myVersion}{version 1.0\xspace}

% ------------------------------------------------------------
%   COMANDOS ÚTILES
% ------------------------------------------------------------

\newcommand{\ie}{i.\,e.}
\newcommand{\Ie}{I.\,e.}
\newcommand{\eg}{e.\,g.}
\newcommand{\Eg}{E.\,g.}

\newcounter{dummy} % Necesario para hipervínculos correctos (índice, bibliografía, etc.)

% ------------------------------------------------------------
%   PAQUETES BÁSICOS
% ------------------------------------------------------------

\usepackage[utf8]{inputenc}
\usepackage[T1]{fontenc}
\usepackage{xspace}

\PassOptionsToPackage{spanish,es-lcroman,es-tabla}{babel}
\usepackage{babel}

\PassOptionsToPackage{square,numbers}{natbib}
\usepackage{natbib}

\PassOptionsToPackage{fleqn}{amsmath}
\usepackage{amsmath}

\PassOptionsToPackage{pdftex}{graphicx}
\usepackage{graphicx}

% Comando para limpiar hasta página izquierda
\newcommand*\cleartoleftpage{%
  \clearpage
  \ifodd\value{page}\hbox{}\newpage\fi
}

% ------------------------------------------------------------
%   TABLAS, FIGURAS Y CAPTIONS
% ------------------------------------------------------------

\usepackage{tabularx}
\setlength{\extrarowheight}{3pt}
\newcommand{\tableheadline}[1]{\multicolumn{1}{c}{\spacedlowsmallcaps{#1}}}
\newcommand{\myfloatalign}{\centering}
\usepackage{caption}
\captionsetup{format=hang,font=small}
\usepackage{subfig}
\usepackage{floatrow}
\floatsetup[table]{capposition=top}

% ------------------------------------------------------------
%   LISTINGS (CÓDIGO FUENTE)
% ------------------------------------------------------------

\usepackage{listings}
\lstset{
  language=[LaTeX]Tex,
  keywordstyle=\color{RoyalBlue},
  basicstyle=\small\ttfamily,
  commentstyle=\color{Green}\ttfamily,
  stringstyle=\rmfamily,
  numbers=left,
  numberstyle=\scriptsize,
  stepnumber=5,
  numbersep=8pt,
  showstringspaces=false,
  breaklines=true,
  frame=single,
  belowcaptionskip=.75\baselineskip
}

% ------------------------------------------------------------
%   HIPERVÍNCULOS
% ------------------------------------------------------------

\PassOptionsToPackage{pdftex,hyperfootnotes=false,pdfpagelabels,colorlinks=true,allcolors=black}{hyperref}
\usepackage{hyperref}



\pdfcompresslevel=9
\pdfadjustspacing=1

\hypersetup{
  colorlinks=true,
  pdfstartpage=1,
  pdfstartview=FitV,
  breaklinks=true,
  pdfpagemode=UseOutlines,
  pageanchor=true,
  plainpages=false,
  bookmarksnumbered,
  bookmarksopen=true,
  bookmarksopenlevel=1,
  hypertexnames=true,
  pdfhighlight=/O,
  urlcolor=webbrown,
  linkcolor=blue,
  citecolor=black,
  linktoc=all,
  % Metadatos PDF
  pdftitle={\myTitle},
  pdfauthor={\textcopyright\ \myName, \myUni},
  pdfcreator={pdfLaTeX},
  pdfproducer={LaTeX with hyperref and classicthesis}
}


\PassOptionsToPackage{smaller}{acronym}
\usepackage{acronym}
\providecommand{\bflabel}{}
\renewcommand{\bflabel}[1]{{#1}\hfill}


% ------------------------------------------------------------
%   AUTORREFERENCIAS (español)
% ------------------------------------------------------------

\makeatletter
\@ifpackageloaded{babel}{
  \addto\extrasspanish{
    \renewcommand*{\figureautorefname}{Figura}
    \renewcommand*{\tableautorefname}{Tabla}
    \renewcommand*{\partautorefname}{Parte}
    \renewcommand*{\chapterautorefname}{Capítulo}
    \renewcommand*{\sectionautorefname}{Sección}
    \renewcommand*{\subsectionautorefname}{Sección}
    \renewcommand*{\subsubsectionautorefname}{Sección}
  }
  \providecommand{\subfigureautorefname}{\figureautorefname}
}{\relax}
\makeatother

% ------------------------------------------------------------
%   CLASSICTHESIS Y PAQUETES ADICIONALES
% ------------------------------------------------------------

\usepackage{classicthesis}
\usepackage[textwidth=25mm]{todonotes}
%\usepackage[textwidth=25mm,disable]{todonotes}
\usepackage{subfiles}
\usepackage{soulutf8}

\newcommand{\todosolve}[2][]{\todo[color=green!50,size=\tiny,#1]{#2}}
\newcommand{\todohint}[2][]{\todo[color=yellow!50,size=\tiny,#1]{#2}}

\makeatletter
\if@todonotes@disabled
  \newcommand{\hlfix}[2]{#1}
\else
  \newcommand{\hlfix}[2]{\texthl{#1}\todo{#2}}
\fi
\makeatother

\usepackage{pdfpages}
\usepackage{lscape}
\usepackage{multirow}
\usepackage{siunitx}
\usepackage{booktabs}
\usepackage{hyphenat}
\usepackage{placeins}
\usepackage[numbered,framed]{matlab-prettifier}
\usepackage{colortbl}
\definecolor{LightCyan}{rgb}{0,0.5,1}
\usepackage[official]{eurosym}
\usepackage{bigstrut}
\usepackage{adjustbox}
\usepackage{gensymb}
\usepackage[margin=2.5cm,top=2.5cm,bottom=2.5cm]{geometry}
\usepackage{makecell}

\usepackage[most]{tcolorbox}

\tcbuselibrary{skins, breakable}
 
\definecolor{reqblue}{HTML}{156082}
 
\newtcolorbox{requisito}[2]{
    enhanced,
    breakable,
    colback       = white,
    colframe      = reqblue,
    colbacktitle  = reqblue,
    coltitle      = white,
    fonttitle     = \bfseries\small,
    title         = {#1 \hfill #2},
    left          = 4pt,
    right         = 4pt,
    top           = 2pt,
    bottom        = 2pt,
    boxrule       = 0.8pt,
    arc           = 2pt,
} 

\usepackage{longtable}
\usepackage{svg}
\usepackage{placeins}
